% Copyright (c) 2026 Eugene Joseph M. Ragasa.
% SPDX-License-Identifier: Apache-2.0
% Author: Eugene Joseph M. Ragasa
% Formal proof completed: 2026-10-20
\documentclass[11pt]{article}
\usepackage[T1]{fontenc}
\usepackage{amsmath,amssymb,booktabs,tabularx}
\usepackage{listings}
\usepackage[margin=1in]{geometry}
\usepackage{xurl}
\usepackage[hidelinks]{hyperref}
\setlength{\emergencystretch}{3em}
\lstset{basicstyle=\ttfamily\small,breaklines=true,columns=fullflexible}
\title{Proof Dossier: Injectivity of Laurent Evaluation on the Unit Circle}
\author{Eugene Joseph M. Ragasa}
\date{Formal proof completed: 2026-10-20}
\begin{document}\sloppy\maketitle
\section{Identity and status}
\begin{tabularx}{\textwidth}{@{}lX@{}}\toprule
Statement ID & \path{BG-PROJECTION-THM-003} \\
Status & proved; no admissions \\
Lean declaration & \path{ProjectKoios.BandGauge.Projection.evalLaurentOnUnitCircle_injective} \\
Owning module & \path{ProjectKoios/BandGauge/Projection/SpectralProjectionRange.lean} \\
Axioms & \texttt{propext}, \texttt{Classical.choice}, \texttt{Quot.sound} \\
\bottomrule\end{tabularx}

\section{Exact mathematical statement}
For $f,g\in\mathbb C[z,z^{-1}]$, if
$f(e^{ik})=g(e^{ik})$ for every $k\in\mathbb R$, then $f=g$.

\section{Source provenance and relationship}
This theorem \emph{formalizes} the uniqueness step used in
\texttt{ksdft2effmass} commit
\texttt{99a46cb59d0c6f27ada7d709f194996ef9619d9f},
\path{docs/publications/papers/Ragasa2026b/v2/appendix_proofs.tex:16--18},
SHA-256
\texttt{614fad3426a941a9a1c5ad4c73ca917b55007b525a1a3a76970d5d7df5af7d19}.
The cited bytes are clean bytes at that immutable commit.

\section{Assumptions, conventions, and delta}
The source writes $z=e^{ik}$ informally. Lean freezes this as
\texttt{LaurentPolynomial $\mathbb C$} evaluated with
\texttt{LaurentPolynomial.eval2} at the complex unit $e^{ik}$. Equality is required for
all real angles. The theorem is a general Laurent uniqueness lemma rather than a claim
specific to the two-band coefficients; this is reusable supporting mathematics for the
exact source step.

\section{Statement-specific prior art}
\begin{tabularx}{\textwidth}{@{}lXXX@{}}\toprule
Relation & Identity & Locator & Role \\
\midrule
formalizes & owner source above & Appendix proof lines 16--18 & required implication in the paper proof \\
background-only & canonical external identity unresolved & unverified & standard polynomial root finiteness; no novelty claim \\
\bottomrule\end{tabularx}

\section{Motivation and contribution boundary}
The theorem prevents a function-level equality on the circle from being silently treated
as coefficient equality. Its contribution is a checked bridge between the source Fourier
language and the Laurent ring, not a claim of originality for polynomial uniqueness.

\section{Proof schematic and argument}
Apply the claim to $q=f-g$. Mathlib's Laurent localization theorem supplies
$n\in\mathbb N$ and $p\in\mathbb C[X]$ with
$p=qz^n$. If $q$ vanishes on the unit circle, then $p$ vanishes there as well. The arc
$\{e^{ik}:0\le k\le\pi\}$ is infinite because the real interval is infinite and the
circle exponential is injective on it. Polynomial root finiteness forces $p=0$.
Since $z^n$ is a unit in the Laurent ring, $q=0$, hence $f=g$.

\section{Lean realization}
The proof uses \path{LaurentPolynomial.exists_T_pow},
\path{Polynomial.eq_zero_of_infinite_isRoot}, \path{Set.Icc_infinite},
\path{Circle.exp_injOn_Icc}, and \path{LaurentPolynomial.isUnit_T}. The toolchain is
Lean 4.35.0 with Mathlib revision
\texttt{56a005855ea17e0b0958f73ec801ae140b0a1679}; no additional dependency is adopted.

\section{Trusted base and axiom report}
\begin{lstlisting}
#print axioms ProjectKoios.BandGauge.Projection.evalLaurentOnUnitCircle_injective
-- [propext, Classical.choice, Quot.sound]
\end{lstlisting}

\section{AI, automation, and comprehension record}
The statement selection, Lean proof search, implementation, and dossier prose were
AI-assisted. Lean elaboration and kernel checking validate the exact declaration. The
maintainer reviewed the reduction to ordinary polynomial roots and the infinite-arc
argument. No independent expert review has occurred.

\section{Scientific and evidence boundaries}
This is exact algebra and elementary complex analysis. It consumes no numerical or
runtime evidence and proves no localization rate, material claim, or universal statement
about spectral projection.

\section{Unresolved questions}
Canonical external prior-art identity remains unresolved through References. A broader
multidimensional torus uniqueness API is not part of this theorem.

\section{Review and revision history}
\begin{tabularx}{\textwidth}{@{}lllX@{}}\toprule
Date & Revision & Process & Outcome \\
\midrule
2026-10-20 & initial & AI-assisted implementation; Lean kernel check & exact unit-circle Laurent uniqueness proved; no independent expert review \\
\bottomrule\end{tabularx}
\end{document}
